\documentclass[11pt,a4paper]{article}

\usepackage[utf8]{inputenc}
\usepackage[T1]{fontenc}
\usepackage[english]{babel}
\usepackage[a4paper,margin=2.4cm]{geometry}
\usepackage{amsmath,amssymb}
\usepackage{bm}
\usepackage{booktabs}
\usepackage{enumitem}
\usepackage{xcolor}
\usepackage{hyperref}
\usepackage{tcolorbox}

\hypersetup{colorlinks=true, linkcolor=blue!50!black, urlcolor=blue!60!black}

\newcommand{\BHF}{B_{\mathrm{hf}}}
\newcommand{\dEQ}{\Delta E_{Q}}
\newcommand{\code}[1]{\texttt{\small #1}}
\DeclareMathOperator*{\argmin}{arg\,min}

\newtcolorbox{nota}[1][Note]{colback=blue!4, colframe=blue!45!black,
  fonttitle=\bfseries, title={#1}, left=4pt, right=4pt, top=3pt, bottom=3pt}
\newtcolorbox{aviso}[1][When to use it]{colback=orange!7, colframe=orange!70!black,
  fonttitle=\bfseries, title={#1}, left=4pt, right=4pt, top=3pt, bottom=3pt}

\title{\textbf{Thickness Correction in \emph{Mossbauer}}\\[3pt]
\large Thick-absorber model: benefits, applicability and usage\\
\normalsize\textcolor{black!60}{(experimental branch \texttt{claude/experimentos-espesor})}}
\author{Department of Physical Chemistry · UAM}
\date{\today}

\begin{document}
\maketitle

\begin{abstract}\noindent
The standard fitting model treats the sample as a \emph{thin} absorber:
transmission is linear in the absorption of each component. In thick
or strongly absorbing samples this approximation fails — the deep lines
\emph{saturate} — and biases depths, linewidths, and especially the areas
(phase fractions).
The experimental branch adds a \emph{thickness correction} via an exponential
saturation model with a single physical parameter $C$. This document explains
the benefits it provides, when it should be applied, how it is applied (what
changes in the fit and in the display), and how it behaves with multiple
subspectra.
\end{abstract}

\tableofcontents
\bigskip

%======================================================================
\section{The problem: thin vs. thick absorber}

The historical model is linear in absorption:
\begin{equation}
  T_{\text{thin}}(v) \;=\; b + s\,v \;-\; A_{\text{tot}}(v),
  \qquad A_{\text{tot}}(v)=\sum_{c} A_c(v),
  \label{eq:thin}
\end{equation}
with $b$ the baseline, $s$ the slope, and $A_c$ the absorption of each component
(sextet/doublet/singlet). It is exact only in the limit of an infinitely thin
absorber. In a real sample of finite thickness, the fraction of resonant photons
reaching the detector decays \emph{exponentially} with the amount of resonant
material: the deepest lines absorb proportionally less than \eqref{eq:thin}
predicts (they \emph{saturate}) and, moreover, they \emph{broaden}.

\begin{nota}
The thickness saturation effect is not a numerical artefact: it is physics of
the transmission experiment. Ignoring it in thick samples causes the fit to
``compensate'' spuriously by inflating linewidths and unbalancing areas.
\end{nota}

%======================================================================
\section{Exponential saturation model}

The correction replaces \eqref{eq:thin} with
\begin{equation}
  \boxed{\;
  T(v) \;=\; b + s\,v \;-\; C\,\Bigl(1 - e^{-A_{\text{tot}}(v)/C}\Bigr)
  \;}
  \label{eq:thick}
\end{equation}
where $C>0$ is the \emph{saturation scale} (maximum achievable absorption amplitude,
$\approx f_s\,b$, with $f_s$ the effective recoil-free fraction of the source).
Its two limits make it interpretable and backward-compatible:
\begin{itemize}[leftmargin=*]
  \item \textbf{Thin limit} $C\to\infty$: since $C(1-e^{-x/C})\to x$,
        \eqref{eq:thin} is recovered exactly. That is, the thin model is the
        special case of zero thickness.
  \item \textbf{Saturation} finite $C$: the effective absorption
        $A_{\text{eff}}=C(1-e^{-A_{\text{tot}}/C})$ is bounded by $C$ and
        compresses the deep lines more as $A_{\text{tot}}/C$ increases.
\end{itemize}
Implemented in \code{core/physics.py} (\code{total\_model}, parameter
\code{sat\_scale}=$C$).

%======================================================================
\section{Benefits}

\begin{enumerate}[leftmargin=*]
  \item \textbf{Unbiased areas and phase fractions.} Saturation suppresses the
        absorption of majority phases (deep lines) more than that of minority
        phases. Without correction, the area ratio is biased in favour of
        weak phases. After \emph{de-saturation}, the recovered areas correspond
        to those of the equivalent thin absorber, which are proportional to
        the abundances (given equal Lamb--M\"{o}ssbauer $f$-factors).
  \item \textbf{Correct depths.} The depth parameter no longer reads the
        apparent (clipped) absorption but represents the real absorption.
  \item \textbf{Less inflated linewidths.} Thickness broadening is no longer
        artificially absorbed into $\gamma$; the fitted linewidths approach
        the intrinsic ones.
  \item \textbf{Better $\chi^2$ for deep spectra.} Where the thin model leaves
        systematic structure in the residual (at the minima), saturation
        removes it without introducing spurious components.
\end{enumerate}

\begin{nota}[Verification]
On a synthetic spectrum saturated with $A_{\max}/C=2$, the fit with thickness
correction reconstructed the spectrum $\sim27\%$ better (RMS) than the thin
model, and recovered the starting scale $C$ ($C=0.41$ versus the true $0.40$).
Under mild saturation it recovered $\sim9\%$ more area (de-saturation).
\end{nota}

%======================================================================
\section{When to apply it}

\begin{aviso}
The thickness correction \textbf{helps} when the sample is thick / strongly
absorbing, and \textbf{is unnecessary} (and should not be forced) when it is thin.
\end{aviso}

Signs that it is worth activating:
\begin{itemize}[leftmargin=*]
  \item Deep absorption: the transmission minimum drops well below the baseline
        (typically $\gtrsim 15$--$20\%$ absorption), or effective thickness
        $t_a\gtrsim 1$.
  \item The thin fit leaves \emph{structured residuals at the minima} (the model
        does not reach the bottom of the deep lines or makes them too wide).
  \item \emph{Reliable phase quantification} (area ratios) is required, not
        just positions.
\end{itemize}
When \emph{not} to apply it:
\begin{itemize}[leftmargin=*]
  \item Thin spectra (absorption of a few \%): the linear model is already exact
        and $C$ becomes undetermined (tends to large values). Keep thin mode.
  \item As a diagnostic: if fitting $C$ converges to a large value and $\chi^2$
        does not improve, the sample is effectively thin.
\end{itemize}

%======================================================================
\section{How it is applied in the fit}

\subsection{Selection and parameter}
In the calibration panel there is a selector \emph{Absorber: \code{thin} /
\code{thickness}} and a saturation scale \emph{slider} $C$
(\code{sat\_scale}), active only in thick mode. $C$ is a global parameter
(like baseline or slope), fittable or fixable.

\subsection{Discrete mode}
The non-linear optimiser (TRF) already fits an arbitrary function, so it suffices
to have the forward model use \eqref{eq:thick}: $C$ enters as one more free
parameter and is refined together with the rest. The algorithm does not change,
only the forward model.

\subsection{Distribution mode \texorpdfstring{$P(\BHF)$}{P(BHF)}}
Here the exponential \emph{breaks the linearity} in the weights $P$ on which
the entire regularised solver relies. The solution is an \emph{inverse transform}
that restores linearity: given the datum $y$, the saturation is inverted to
obtain the observed linear absorption
\begin{equation}
  A_{\text{obs}}(v) \;=\; -\,C\,\ln\!\Bigl(1 - \tfrac{b + s\,v - y(v)}{C}\Bigr),
  \label{eq:invtransform}
\end{equation}
which \emph{is} linear in $P$ ($A_{\text{obs}}=K\,P+\dots$). The usual
regularised linear system (Tikhonov/TV, non-negativity) is solved on
$A_{\text{obs}}$, and the model is \emph{re-saturated} with \eqref{eq:thick} to
compute residuals and the displayed curve. The noise propagates through the
transform,
\begin{equation}
  \sigma_{A} \;=\; \sigma_y\,\frac{dA_{\text{obs}}}{dy} \;=\; \sigma_y\,e^{A_{\text{obs}}/C},
\end{equation}
so that saturated channels are correctly weighted. An outer non-linear loop
(separable VARPRO scheme) refines $(b,s,C)$ while the inner linear solve
recovers $P$ at each iteration.

\begin{nota}
The transform \eqref{eq:invtransform} requires $b+s\,v-y < C$ in every channel
(the apparent absorption cannot exceed the saturation plateau); the program clips
the argument to the valid domain $(0,1]$ before the logarithm.
\end{nota}

%======================================================================
\section{What changes in the view}

\begin{itemize}[leftmargin=*]
  \item The \textbf{model curve} displayed is the \emph{saturated} transmission
        \eqref{eq:thick}: it is what is compared with the data and what generates
        the residual.
  \item The \textbf{depths/areas} shown and exported are those of the equivalent
        thin absorber (de-saturated): therefore, when the correction is activated,
        relative areas may change with respect to the thin fit — that is precisely
        the intended correction.
  \item The fitted value of $C$ (saturation scale) appears among the global
        parameters. A large $C$ indicates a practically thin sample.
  \item In distribution mode, $P(\BHF)$ is displayed already de-saturated; its
        integrated area reflects the real fraction.
\end{itemize}

%======================================================================
\section{Multiple subspectra}

Saturation is a property of the \emph{complete absorber}, not of each phase
separately: photons at a given velocity traverse all the material and ``see''
the sum of all absorptions at that velocity. Therefore \eqref{eq:thick} saturates
the \textbf{total absorption} $A_{\text{tot}}=\sum_c A_c$ with a \emph{single}
common scale $C$ for all subspectra, not each component separately:
\begin{equation}
  T(v) = b + s\,v - C\Bigl(1 - \exp\!\bigl[-\tfrac{1}{C}\textstyle\sum_c A_c(v)\bigr]\Bigr).
\end{equation}
Practical consequences:
\begin{itemize}[leftmargin=*]
  \item \textbf{Correct physical coupling.} Where two subspectra overlap, the
        saturation is greater (the local absorption is the sum); the model
        captures this automatically.
  \item \textbf{Phase quantification.} After de-saturation, the area ratios
        between subspectra recover their unbiased value, which is the goal in
        phase mixtures (e.g. multiple Fe environments).
  \item \textbf{A single $C$.} No scale is fitted per component (that would be
        non-physical and would degenerate with the depths): $C$ is global and
        describes the effective thickness of the sample.
  \item \textbf{Compatibility.} In distribution mode, added ``sharp'' sextets
        and the distribution itself share the same correction, because all
        contribute to $A_{\text{tot}}$.
\end{itemize}

\begin{nota}[Model limitation]
This is an effective exponential saturation (Lambert--Beer), which captures the
dominant thickness effect (compression, broadening, unbiased areas) with a single
parameter. It is not the full Margulies--Ehrman transmission integral (convolution
with the source line), which would be a second-order refinement rarely needed in
practice.
\end{nota}

%======================================================================
\section{Summary}

\begin{center}
\begin{tabular}{@{}ll@{}}
\toprule
\textbf{Aspect} & \textbf{Thickness correction}\\
\midrule
Model & $T=b+s v - C(1-e^{-A_{\text{tot}}/C})$ \\
New parameter & $C$ (saturation scale), global, $C\to\infty$ = thin\\
When & thick samples / deep absorption / phase quantification\\
Discrete & $C$ free in TRF; only the forward model changes\\
Distribution & inverse transform $A_{\text{obs}}=-C\ln(1-(b+sv-y)/C)$ + VARPRO $(b,s,C)$\\
View & saturated curve; de-saturated areas/$P$; $C$ shown\\
Subspectra & single $C$ saturates the sum $\sum_c A_c$ (physical coupling)\\
Key benefit & unbiased areas/phase fractions, better linewidths and $\chi^2$\\
\bottomrule
\end{tabular}
\end{center}

\end{document}
